\documentclass[11pt,a4paper]{article}
\usepackage[margin=1in]{geometry}
\usepackage[T1]{fontenc}
\usepackage[utf8]{inputenc}
\usepackage{lmodern}
\usepackage{microtype}
\usepackage{hyperref}
\usepackage{booktabs}
\usepackage{graphicx}
\usepackage{enumitem}
\setlist{nosep}

\hypersetup{
  colorlinks=true,
  linkcolor=blue,
  urlcolor=blue,
  citecolor=blue,
  pdftitle={GSoC 2026 Task 2b CMS Jet Super-Resolution},
  pdfauthor={CMS\_E2E}
}

\title{ML4SCI GSoC 2026 --- Task 2b\\\large Super-Resolution at the CMS Detector (Jet Images)}
\author{Repository \texttt{KenWuqianghao/CMS\_E2E}}
\date{\today}

\begin{document}
\maketitle
\tableofcontents
\vspace{1em}
\hrule
\vspace{1em}

\section{Objective}
Train a \textbf{Generative Adversarial Network} to map low-resolution calorimeter jet images to high-resolution targets, using the official ML4SCI parquet release (CERNBox link in the GSoC PDF). The submission documents modeling and optimization choices as required by the task statement.

\section{Dataset and labels}
\begin{itemize}
  \item \textbf{Low resolution:} \texttt{X\_jets\_LR}, shape $3 \times 64 \times 64$.
  \item \textbf{High resolution:} \texttt{X\_jets}, shape $3 \times 125 \times 125$.
  \item \textbf{Class label:} \texttt{y} $\in \{0,1\}$ (quark vs.\ gluon), stratified in train/validation/test splits (default $80\%/10\%/10\%$).
\end{itemize}
The public shards use a \textbf{single Parquet row group per file}; the loader therefore \textbf{streams record batches} and materializes only the row indices needed for training, avoiding full-file loads per sample.

\section{Model}
\begin{itemize}
  \item \textbf{Generator (SRGAN-style):} bicubic upsampling of LR to $125\times 125$ as an explicit baseline, plus a residual CNN (residual blocks, InstanceNorm) predicting a correction added to that baseline.
  \item \textbf{Discriminator (PatchGAN):} convolutional stack with optional \textbf{class conditioning} via embedding-modulated feature maps (quark/gluon).
  \item \textbf{Losses:} least-squares GAN objectives; \textbf{L1} reconstruction on normalized HR; small weight on \textbf{adversarial} loss; optional \textbf{energy-matching} term on denormalized total deposited energy ($\sum \max(0,x)$ over channels and pixels).
\end{itemize}

\section{Training and protocols}
Two tiers are supported in code and documentation:
\begin{enumerate}
  \item \textbf{Host-constrained (reference):} materialized subsets via \texttt{--max-train-samples} / \texttt{--max-val-samples}, fewer epochs, optional \texttt{--memmap-train} when RAM is insufficient for full tensors.
  \item \textbf{Full-scale:} no subset caps, \texttt{--epochs 20}, \texttt{--memmap-train} as needed for $\sim$full training index sets.
\end{enumerate}
Validation logs \texttt{val\_l1}, bicubic baseline \texttt{val\_bicubic\_l1}, and \texttt{val\_energy\_mae} to \texttt{history.csv}. Checkpoints: \texttt{checkpoint\_best.pt} (best val L1), \texttt{checkpoint\_last.pt}.

\section{Evaluation}
Script \texttt{super\_resolution/evaluate.py} reports JSON metrics on the test split (full test by default), including SR vs.\ HR and a \textbf{bicubic} baseline: PSNR/SSIM, mean total-energy ratio, energy MAE, mean radial-profile L1, and a toy linear probe on per-image channel means. Results are written as structured JSON (e.g.\ \texttt{metrics\_report.json}).

\section{Reference results (checked in)}
Using the \textbf{full\_run} configuration (subset caps, 8 epochs, $\lambda_{\mathrm{energy}}=0.05$) on the full test split ($n=13\,932$), illustrative numbers from \texttt{super\_resolution/runs/full\_run/metrics\_report.json}:
\begin{center}
\begin{tabular}{@{}lrr@{}}
\toprule
Metric & SR & Bicubic \\
\midrule
PSNR mean (dB) & 44.28 & 44.28 \\
SSIM mean & 0.9935 & 0.9934 \\
Energy MAE vs.\ HR & 1.20 & 21.29 \\
Mean radial profile L1 vs.\ HR & 0.00040 & 0.00132 \\
\bottomrule
\end{tabular}
\end{center}
The energy penalty materially improves \textbf{energy fidelity} vs.\ bicubic; image PSNR/SSIM remain close to the trivial resize because the task geometry is dominated by a smooth upscale.

An \textbf{ablation} with \texttt{--lambda-energy 0} (\texttt{ablation\_no\_energy/}) shows degraded energy scale (e.g.\ energy MAE $\approx 153$ vs.\ HR) despite similar PSNR, illustrating the need for physics-aware terms beyond pixel loss.

\section{Code layout}
\texttt{super\_resolution/data.py}, \texttt{models.py}, \texttt{train.py}, \texttt{evaluate.py}; narrative in \texttt{MODEL\_DISCUSSION.md} and \texttt{super\_resolution/README.md}.

\section*{Report artifact}
This PDF is built from \texttt{docs/GSoC26\_Task2b\_Report.tex} in \texttt{CMS\_E2E}.

\end{document}
